% =============================================================================
% CHAPTER 6: Reflection on Learning
% =============================================================================
% SPDX-FileCopyrightText: 2026 Frank Langbein <frank@langbein.org>
% SPDX-License-Identifier: LPPL-1.3c
%
% Suggested sections: Critical narrative of the experience; Analysis and application;
% Implications for future practice. Focus on transferable learning and (where
% applicable) ``double loop learning'': impact on assumptions and concepts, not
% only skills. Optional alternatives: ``What I would do differently''; ``Transferable
% skills and concepts''. Omit this chapter if not required by your institution;
% comment out the \chapter and \input in main.tex.
% =============================================================================

\newthought{Committing to a project} of this caliber and the time to focus a majority of the Author's time to it has provided some critical reflections for the Author about the project as a whole and how they learn independently.

\section{Critical narrative of the experience}\label{sec:challenges}
% Key challenges, turning points, and how they affected your understanding.
Producing a solution tailored for an end user that is not the Author was a learning experience. Most of the applications that the Author created beforehand have been either coursework assignments or simple programs that the Author could always amend. This encouraged the Author to produce as cohesive and as self-sufficient of a product as possible. The Author encountered a variety of challenges and would like to highlight some of them and what they learnt from them.
\begin{itemize}
    \item \textbf{Working with Embedded Systems} - Specialised hardware is designed for a specific purpose and has specific requirements for it to work alongside the system as a whole. The Author acknowledged that many of these requirements may not be obvious at first so prior research into the embedded systems and their components is vital. This can be seen in the OS incompatibility with the GrovePi+ shield for the Raspberry Pi. The surface level information the Author was provided with was insufficient as it was the further research that made them realise the kernel incompatibility and deprecated URLs for certain packages. This highlighted the importance of research that cannot be overstated, which the Author will take into consideration from this point forward.  
    \item \textbf{Planning for the Long-Term} - Working with different parties and stakeholders within a project naturally creates trouble in terms of scheduling, which the Author did not fully grasp. The Author left it too late for any ethical application to be sent off in time and then arrange for user testing, which should have happened during the beginning of the project to ensure the possibility of user testing. This can also be reflected in more continuous contact with the project supervisor. Had the Author engaged with the supervisor more, they could have potentially identified the need for ethical approval for user testing much sooner and reached out to more researchers to obtain their feedback on the proposed solution, helping direct development with clear tailored objectives for the proposed solution, rather than what the Author assumed would be suitable.
\end{itemize}

Ultimately, the Author has learnt a lot about how a development pipeline with weekly meetings would work. The Author now knows how to appropriately plan for potential hardship within a project and when to ask for help when it is necessary. The Author learnt how to do more targeted research and how to find basis or possible disputes in the assumptions that the Author may have previously had. Finally, the Author learnt how to engage with the public on a piece of research and how to communicate effectively with their colleagues to boost the productivity of their research into a solution.

